\documentclass[11pt]{article}
\usepackage[T1]{fontenc}
\usepackage{lmodern,microtype,amsmath,amssymb,amsthm,booktabs,array}
\usepackage[margin=1in]{geometry}
\usepackage[hidelinks]{hyperref}
\hypersetup{pdftitle={Acyclicity and simultaneous protection of individual core maxima},pdfauthor={Maxwell Flanagan},pdfsubject={International Journal of Game Theory submission preparation}}
\usepackage[round,authoryear]{natbib}
\usepackage{enumitem}
\setlist{nosep}
\newtheorem{theorem}{Theorem}[section]
\newtheorem{lemma}[theorem]{Lemma}
\newtheorem{proposition}[theorem]{Proposition}
\newtheorem{corollary}[theorem]{Corollary}
\theoremstyle{remark}\newtheorem{remark}[theorem]{Remark}
\newcommand{\core}{\mathcal C}
\newcommand{\lex}{\operatorname{leximin}}
\newcommand{\R}{\mathbb R}
\newcommand{\N}{\mathbb N}
\title{Acyclicity and simultaneous protection\\of individual core maxima}
\author{Maxwell Flanagan\\\small New York City, United States\\\small\href{mailto:maxwell.flanagan@gmail.com}{maxwell.flanagan@gmail.com}}
\date{October 17, 2026}
\begin{document}
\maketitle
\begin{abstract}
A stable allocation need not protect each participant relative to that
participant's most favorable stable payoff. We characterize when sector
interactions permit a uniform simultaneous guarantee of these individual core
maxima. Fix a connected graph on $K\ge2$ equal-sized sectors. Full teams contain
one player from every sector, their values sum nonnegative pair values on the
allowed edges, and only full teams produce. Coalition worth is the maximum
value of an integral matching. Among markets with positive worth and nonempty
core, both full leximin and the best choice of any core allocation have sharp
size-independent guarantee $1/K$ on trees and zero on graphs containing a
cycle. The tree proof describes the entire core by cut coordinates and uses
clipping at every lexicographic level. For a triangular family with pair values
in $\{0,1,2\}$, the entire core is a simplex and the optimal simultaneous
fraction is exactly $1/(2q+3)$. Selective subdivision transfers vanishing
protection to every cyclic allowed graph while preserving a positive core
allocation. The cyclic obstruction therefore concerns what stability can
achieve, not merely the choice of a fairness rule. The classification is
conditional on complete-team production and permits zero-valued allowed edges.
\end{abstract}
\noindent\textbf{Keywords:} assignment game; core; leximin; individual ideal;
simultaneous protection; sector graph.
\par\noindent\textbf{JEL classification:} C71; C78.

\section{Introduction}
When production requires partners from several sectors, coalitional stability
and individual protection address different questions. The core rules out a
group earning more by producing separately. It does not ensure that a player
receives a substantial share of the most that player could obtain in another
stable agreement. This distinction matters when the core is used as a feasible
set for bargaining: a payoff may be stable yet small relative to the same
participant's attainable stable payoff.

For an entire core $\core$, write $H_i=\max_{y\in\core}y_i$. We ask whether
one stable allocation can guarantee $p_i\ge\lambda H_i$ simultaneously for all
players, with $\lambda$ depending only on the number and pattern of sector
interactions, not the number of players. We also ask whether full leximin,
which successively maximizes the smallest nominal payoffs, attains such a
uniform guarantee. The benchmark $H_i$ is an attainable core payoff, not an
outside option or a claim that all maxima can be attained together. We keep
this benchmark fixed when comparing allocations within a market.

Our model describes a production project that needs one participant from
every sector. Pair values describe contributions to a complete project's
surplus; an edge does not by itself authorize a productive bilateral coalition.
The same sharp graph boundary governs both questions. Trees guarantee $1/K$,
and full leximin attains that bound for everyone at once. Every cyclic allowed
graph instead admits growing markets with vanishing best simultaneous
protection, even though each market has a strictly positive core allocation.
Thus the negative result cannot be repaired by selecting a different point
of the same core. The poorly protected player may depend on that point;
each individual maximum is nevertheless attained somewhere in the core.

The proof has three parts. A cut-coordinate description of the entire tree
core yields a clipping operation that improves any violation of the $1/K$
floor. A fixed-alphabet triangular construction has an explicit simplex core
and exact optimal fraction $1/(2q+3)$. Finally, selective subdivision transfers
a unit partial-matching shortfall to all cyclic allowed graphs. The latter two
steps keep the number of sectors fixed as market size grows. The results
measure distributional possibilities under stability, not a loss in aggregate
surplus, typical-market performance, or a change in productive efficiency.

\paragraph{Related work.}
The multisided assignment model and its core have a substantial literature,
including \citet{Quint1991} and \citet{Lucas1995}. The strict
supplier--firm--buyer model is associated with \citet{Stuart1997};
\citet{AtaySolymosi2018}, Definition~3 and equations~(1)--(2), provide a
later full formulation with signed adjacent-pair values and zero-truncated
triple totals. Our nonnegative three-sector path is a subclass. We claim
neither the model nor chain-core nonemptiness as new.

The closest graph comparison is \citet{AtayNunez2019}. Their Theorem~1
expresses tree cores through bilateral cores, Corollary~1 gives a tree/cycle
strong-balancedness boundary, and Proposition~2 gives simultaneous maxima
within a sector. Their connected partial-sector teams also produce value.
Our core instead enforces strict full-team production and the cyclic theorem
is conditional on its nonemptiness. Thus neither their core decomposition
nor their empty-core cyclic examples establish the protection classification
here. The elementary comparison below shows explicitly how partial-team
production changes the individual maxima.

Simultaneous ideal protection has a direct bilateral precedent: the
half-optimal core of N\'u\~nez and Rafels requires at least half of each
player's maximum core payoff. The original working text is N\'u\~nez and
Rafels, \emph{Bargained stable allocations in assignment markets}, Barcelona
Economics Working Paper~153, December 28, 2004, Definition~6, printed p.19,
with sharpness in Example~7, p.20
(\url{https://www.bse.eu/sites/default/files/working_paper_pdfs/153.pdf}).
The bilateral bargaining midpoint is also studied by
\citet{NunezRafels2008}. Neither the benchmark nor the bilateral half guarantee
is proposed as a new fairness concept.

\citet{Gangam2026}, p.1621, give a unit-capacity path whose equitable payoff
is $1/r$ on each of its $r$ edges. Its entire core is the unit simplex;
therefore each ideal equals one and even the best simultaneous fraction is
$1/r$. This last conclusion is an immediate inference from that example.
Its productive coalition size grows with $r$, whereas our obstruction already
uses strict triples and pair values in $\{0,1,2\}$. The distinction is fixed
production-team size within a restricted class, not the first occurrence of
vanishing individual protection in cooperative games.

Extreme-core descriptions for assortative multisided markets
\citep{Martinez2020} and simultaneous extrema under supermodularity and local
monotonicity \citep{Martinez2026} impose different restrictions. A general
2023 core-redescription preprint of the same authors also retains original
inequalities outside a central strip: Mart\'inez-de-Alb\'eniz, Rafels and
Ybern, \emph{Extreme core points in multisided assignment problems}, 2023
preprint, Theorem~2 and Theorems~3--4
(\url{https://doi.org/10.2139/ssrn.4369524}).
That general representation may apply to our inputs; non-supermodularity
does not exclude it. It does not itself give the quantitative protection
classification proved here. We do not assert an exclusion under every
relabeling or assume an arbitrary core representation preserves the payoff
labels, selector and pair-value restrictions relevant to this question.

Matching duality, individual ideals, leximin, convex averaging and closure
of difference constraints under minimum are established ingredients. The
proposed contribution is their specific constructive use to classify
protection on the entire strict full-team core, with a branching-tree
argument and a bounded-alphabet obstruction on every cyclic allowed graph.
Full-original comparisons with Quint, Lucas and Stuart remain incomplete
under freely available access; their basic attribution is retained, without
assuming uninspected detailed statements or claiming exhaustive priority
clearance. No theorem below relies on such an uninspected statement.

\section{Model and main classification}
Fix a connected simple graph $G=(V,E)$ with $K=|V|\ge2$. Each sector $v\in V$
contains players $(v,j)$, $1\le j\le m$, for a common finite $m\ge1$.
An allowed edge $uv$ carries a nonnegative matrix $A_{uv}$; zero matrices
are permitted. All full teams $t=(j_v)_{v\in V}$ are available, with worth
\[
 w(t)=\sum_{uv\in E}A_{uv}(j_u,j_v).
\]
Only full teams produce value. For any coalition $S$, its worth $v(S)$ is the
maximum sum of worths of player-disjoint full teams contained in $S$.
In particular every singleton has worth zero. Set $T=v(N)>0$ and require
that the entire core $\core$ be nonempty. Equivalently,
\begin{equation}\label{eq:core}
 \core=\{p\ge0:\ \sum_i p_i=T,\ \sum_{i\in t}p_i\ge w(t)
                    \text{ for every full team }t\}.
\end{equation}
Indeed these conditions imply every coalition inequality by summing over
its integral matching and retaining nonnegative unmatched payoffs.

The compact convex set $\core$ has coordinate maxima $H_i$. Let $x$ be its
full ascending-sorted leximin allocation; all labeled coordinates, including
permanent zeros, remain in the optimization. Existence follows by successive
maximization of order statistics on compact optimal sets. Uniqueness follows
by averaging two distinct purported optima: at their first differing payoff
level the average reduces its multiplicity without changing lower levels.
Define
\[
 c(G)=\inf_M\min_{i:H_i>0}\frac{x_i}{H_i},\qquad
 \rho(M)=\max_{p\in\core}\min_{i:H_i>0}\frac{p_i}{H_i},\qquad
 R(G)=\inf_M\rho(M),
\]
where $M$ ranges over the stated markets on fixed $G$, allowing $m$ to grow.
Thus $c(G)\le R(G)$. The ratios are formed only when $H_i>0$; other
coordinates are not deleted from the core or leximin problem.

\begin{theorem}\label{thm:main}
For every fixed connected allowed graph on $K\ge2$ sectors,
\[
 c(G)=R(G)=
 \begin{cases}1/K,&G\text{ is a tree},\\0,&G\text{ contains a cycle}.
 \end{cases}
\]
For every cyclic $G$ and integer $q\ge1$, the upper witnesses have $3q+1$
players per sector, pair entries in $\{0,1,2\}$, and a strictly positive
core allocation. One target has ideal $1$ and full leximin payoff
$1/[K(q+1)]$, while $\rho(M)\le1/(2q+2)$.
For the source triangle, the exact stronger formula is
$\rho(M_q)=1/(2q+3)$.
\end{theorem}

Every fixed finite market still has $\rho(M)>0$: average one core maximizer
for each of its $n_+$ positive ideals to obtain $p_i\ge H_i/n_+$.
The zero in the theorem is an infimum over sizes, not a finite-market zero.
Also, the theorem concerns allowed edges, not strictly positive matrices on
every allowed edge. A single-sector edge-sum model has no positive-worth
instance and is excluded.

\paragraph{Why complete-team production matters.}
Consider three sectors with one player each on the path $1$--$2$--$3$, with
both edge values equal to one. In the present model only the triple produces,
so $T=2$ and the core is $\{p\ge0:p_1+p_2+p_3=2\}$. All three ideals equal
two and the best simultaneous fraction is $1/3$. If adjacent pairs can also
produce their edge values, the core additionally requires
$p_1+p_2\ge1$ and $p_2+p_3\ge1$. Its ideals become $(1,2,1)$; the point
$(1/2,1,1/2)$ achieves the optimal fraction $1/2$, since the ideals sum to
four while the total remains two. Thus admitting partial teams changes the
benchmark as well as the feasible set. A tree theorem for that model does
not transfer the same relative-payoff guarantee to the present core.


\section{Trees: entire-core coordinates and clipping}
Root the tree at an arbitrary sector $r$. Choose an optimal perfect matching
on each edge and propagate labels from the root so that each is diagonal.
There is no conflicting propagation on a tree. Write $d_{ab}(j)=A_{ab}(j,j)$
for an oriented parent--child edge $a\to b$. Any full-team matching projects
to edge matchings; conversely the diagonal full-team matching attains the sum
of these edge optima. Hence
$T=\sum_j\sum_{ab\in E}d_{ab}(j)$.
Nonnegative weights and equal sizes permit completion of partial matchings
without loss. Summing optimal nonnegative bilateral dual covers over incident
edges gives a full-team cover of total $T$, so tree cores are nonempty.
Every diagonal team is tight throughout the core.

For a nonroot vertex $b$, let $S_b$ be its descendant subtree, $s_b=|S_b|$,
and define
\[
 U_b^p(j)=\sum_{v\in S_b}p_v(j)
                -\sum_{e\text{ internal to }S_b}d_e(j),\qquad U_r^p(j)=0.
\]
\begin{lemma}\label{lem:cut}
The entire core is in bijection with the cut vectors satisfying
\begin{align}
 U_b(i)-U_b(j)&\ge A_{ab}(j,i)-d_{ab}(j),\label{eq:cut}\\
 p_v(j)&=U_v(j)-\sum_{b\text{ child of }v}U_b(j)
                   +\sum_{b\text{ child of }v}d_{vb}(j)\ge0.\label{eq:reconstruct}
\end{align}
\end{lemma}
\begin{proof}
For necessity, take label $i$ on $S_b$ and label $j$ on its complement in a
full team, and subtract the tight diagonal team at $j$. The single cut edge
is $ab$; internal diagonal terms cancel to give~\eqref{eq:cut}.
Subtracting child-subtree sums gives~\eqref{eq:reconstruct}; at the root use
diagonal tightness. Conversely reconstruction telescopes at each fixed label
to the diagonal team value and therefore gives total $T$. At arbitrary labels,
\[
 \sum_v p_v(i_v)=\sum_{a\to b}
 [U_b(i_b)-U_b(i_a)+d_{ab}(i_a)]
 \ge\sum_{a\to b}A_{ab}(i_a,i_b).
\]
Thus every full team is covered. Nonnegativity and~\eqref{eq:core} imply core
membership. Descendant sums recover the cut coordinates, proving bijectivity.
\end{proof}

\begin{lemma}[Clipping]\label{lem:clip}
For $x,y\in\core$ and $\beta\ge0$, set
$Z_b(j)=\min\{U_b^x(j),U_b^y(j)+s_b\beta\}$ for $b\ne r$, and $Z_r=0$.
The reconstructed $z$ belongs to $\core$ and satisfies
\[
 z_v(j)\ge\min\{x_v(j),y_v(j)+\beta\}\quad(v\ne r),\qquad
 z_r(j)\ge\max\{x_r(j),y_r(j)-(K-1)\beta\}.
\]
\end{lemma}
\begin{proof}
Difference inequalities $u(i)-u(j)\ge c_{ij}$ are closed under uniform
translation and coordinatewise minimum: both candidate values at $i$ are at
least their respective value at $j$ plus $c_{ij}$. Thus $Z$ satisfies all
structural inequalities. At a nonroot vertex, if its selected cut is the $x$
cut, each child cut is at most its $x$ cut and reconstruction gives $z_v\ge x_v$.
If it is the shifted $y$ cut, use $s_v=1+\sum_b s_b$ to obtain
$z_v\ge y_v+\beta$. These cases include ties. At the root, upper-bound each
child cut first by its $x$ value and then by its shifted $y$ value; the child
subtrees have total size $K-1$. This gives both root inequalities. All
reconstructed payoffs are nonnegative; Lemma~\ref{lem:cut} finishes the proof.
\end{proof}

Fix a target at the root sector and put $\alpha=x_i$ for full leximin $x$.
If $H_i>K\alpha$, choose an ideal-attaining $y$ and
$\alpha<\beta<H_i/K$. Clipping gives
$z_i\ge H_i-(K-1)\beta>\beta>\alpha$.
Every old coordinate at most $\alpha$ weakly increases, and every old
coordinate above $\alpha$ remains strictly above it. Thus, for each
$\gamma\le\alpha$, the number of coordinates below $\gamma$ cannot
increase, and the set of coordinates at most $\alpha$ strictly shrinks
because the target leaves it. Decreases above $\alpha$ do not affect this
comparison. If a lower level
increases, its first changed multiplicity decreases; otherwise the
multiplicity of $\alpha$ decreases. Either case strictly improves the sorted
vector, a contradiction. Consequently $x_i\ge H_i/K$, even when $x_i=0$.
Re-rooting the tree and relabeling its optimal matching change only the
coordinates used in the proof, not the chosen leximin allocation $x$.
The target was arbitrary, so this same $x$ protects every player and
$\rho(M)\ge1/K$ as well.

For sharpness take $m=1$ and edge values totaling one. The sole full team
has worth one and every proper coalition has worth zero. Its core is the
$K$-coordinate unit simplex; $H_i=1$, leximin is uniform, and
$\rho=1/K$. This proves both tree equalities. It also gives the sharp distance
bound $\|x-y\|_1/T\le2(1-1/K)$ for any tree-core point $y$, since
$\min(x_i,y_i)\ge y_i/K$ and both totals are $T$.
\section{A cyclic family with vanishing protection}
Fix an integer $q\geq1$. Use the common team-label set
\[
           \mathcal T=\{L_0\}\cup\{O_s,L_s,R_s:1\leq s\leq q\}.
\]
For every label $t$, create players $a_t,b_t,c_t$ in the respective
sectors. Each sector has $3q+1$ players. Define
$w(a_t,b_u,c_v)=A_{tu}+B_{tv}+C_{uv}$ by the nonzero entries in Table~\ref{tab:entries};
all unspecified pair entries are zero.

\begin{table}[htbp]
\centering
\caption{Nonzero pair values in the triangular market}\label{tab:entries}
\begin{tabular}{lll}\toprule
Pair entries & Range & Value\\\midrule
$C_{L_s,L_s}$ & $0\leq s\leq q$ & $1$\\
$C_{R_s,R_s}$ & $1\leq s\leq q$ & $1$\\
$A_{O_s,O_s},\ B_{O_s,O_s}$ & $1\leq s\leq q$ & $2$\\
$A_{O_s,R_s},\ B_{O_s,L_s}$ & $1\leq s\leq q$ & $2$\\
$C_{R_s,L_s},\ C_{L_{s-1},O_s}$ & $1\leq s\leq q$ & $1$\\\bottomrule
\end{tabular}
\end{table}

There are $8q+1$ nonzero entries. The diagonal perfect matching consists
of $(a_t,b_t,c_t)$ for all $t\in\mathcal T$. Each $L$ or $R$ triple has
value one and each $O$ triple has value four, so this matching is worth
$6q+1$.

\subsection{Two explicit full-core covers}
Put $\tau=1/[3(q+1)]$. Define allocations $u$ and $y$ in $(a,b,c)$ order:
\begin{equation}\label{eq:witnesses}
\begin{array}{c|c|c}
\text{team}&u&y\\\hline
L_s&\bigl(\tau,\ 1-(3s+2)\tau,\ (3s+1)\tau\bigr)&(0,0,1)\\
R_s&\bigl(\tau,\ 1-2\tau,\ \tau\bigr)&(0,1,0)\\
O_s&\bigl(4-(3s-1)\tau,\ \tau,\ (3s-2)\tau\bigr)&(3,0,1).
\end{array}
\end{equation}
The $L$ row uses $0\leq s\leq q$; the other rows use $1\leq s\leq q$.
Every coordinate of $u$ is at least $\tau>0$, and $y$ is nonnegative.
The smallest $L$-sector $b$ coordinate is attained at $s=q$ and equals
$\tau$; the smallest $O$-sector $c$ coordinate is attained at $s=1$
and also equals $\tau$. Each diagonal triple receives its value at
both points, so both allocations have total $6q+1$.

We check every triple cover, including all mixed-team triples. If its
$a$ player has an $L$ or $R$ label, its $A$ and $B$ entries are zero.
Table~\ref{tab:weak-cover} lists the possible positive $C$ pairs and the associated two-player sums.

\begin{table}[htbp]
\centering
\caption{Coverage when the first-sector player has a weak label}\label{tab:weak-cover}
\begin{tabular}{lll}\toprule
$(b,c)$ labels & $u_b+u_c$ & $y_b+y_c$\\\midrule
$(L_s,L_s)$ & $1-\tau$ & $1$\\
$(R_s,R_s)$ & $1-\tau$ & $1$\\
$(R_s,L_s)$ & $1+(3s-1)\tau$ & $2$\\
$(L_{s-1},O_s)$ & $1-\tau$ & $1$\\\bottomrule
\end{tabular}
\end{table}

All four pair values are one. Adding $u_a=\tau$ or $y_a=0$ covers the
triple. Zero $C$ entries are covered by nonnegativity.

Now suppose its $a$ player is $a_{O_s}$. If at most one of its $A$ and
$B$ entries is positive, the triple value is at most three, covered by
$u_{a_{O_s}}\geq3+4\tau$ and $y_{a_{O_s}}=3$.
If both entries are positive, exactly the four cases in Table~\ref{tab:strong-cover} remain.

\begin{table}[htbp]
\centering
\caption{Coverage when both adjacent pair values are positive}\label{tab:strong-cover}
\begin{tabular}{llll}\toprule
$(b,c)$ labels & $w$ & $u_a+u_b+u_c$ & $y_a+y_b+y_c$\\\midrule
$(O_s,O_s)$ & $4$ & $4$ & $4$\\
$(O_s,L_s)$ & $4$ & $4+3\tau$ & $4$\\
$(R_s,O_s)$ & $4$ & $5-3\tau$ & $5$\\
$(R_s,L_s)$ & $5$ & $5$ & $5$\\\bottomrule
\end{tabular}
\end{table}

Since $\tau\leq1/6$, all are covers. The cases also show that every
triple value is at most five.

The matching gives worth at least $6q+1$, while either nonnegative
cover bounds every matching by its total $6q+1$. Hence the true integral
worth is $T=6q+1$ and both $u,y$ belong to the entire core. No fractional
matching value is substituted for coalition worth.

\subsection{A partial matching determines the target}
For each $s=1,\ldots,q$, take the triples
\[
 N_s=(a_{O_s},b_{R_s},c_{L_s}),\qquad
 E_s=(a_{L_{s-1}},b_{L_{s-1}},c_{O_s}).
\]
Their values are five and one, respectively. They form a player-disjoint
partial matching: the used $a$ labels are the $O$ labels and
$L_0,\ldots,L_{q-1}$; the used $b$ labels are the $R$ labels and the same
$L$ labels; the used $c$ labels are $L_1,\ldots,L_q$ and the $O$ labels.
Its total is $6q$. Its unmatched players are exactly
\begin{equation}\label{eq:unmatched}
 U=\{c_{L_0},a_{L_q},b_{L_q}\}
       \cup\{a_{R_s},b_{O_s},c_{R_s}:1\leq s\leq q\},
       \qquad |U|=3q+3.
\end{equation}
Summing the partial-matching inequalities and subtracting from efficiency
gives, at every core point $p$,
\begin{equation}\label{eq:budget}
                       \sum_{a\in U}p_a\leq1.
\end{equation}
Thus a common payoff floor cannot exceed $1/(3q+3)=\tau$. The point
$u$ attains this floor, so the first global leximin level is exactly
$\tau$. At every point of the first-stage optimal set, all coordinates are
at least $\tau$, and $|U|\tau=1$. Inequality~\eqref{eq:budget} therefore
fixes every coordinate of $U$ at $\tau$ on that entire set. In particular,
the full leximin target $c_{L_0}$ is fixed at $\tau$ through every later
refinement. We do not need, and do not assert, that $u$ is the complete
final leximin vector.

The diagonal $L_0$ triple is tight at every core point and has value one.
Nonnegativity implies $H_{c_{L_0}}\leq1$, while the cover $y$ attains
one at that coordinate. Consequently
\[
       H_{c_{L_0}}=1,\qquad x_{c_{L_0}}=\frac{1}{3(q+1)}.
\]
This proves the triangular nominal-leximin conclusion. The next section
characterizes the best simultaneous protection over the entire core.

\section{The entire triangular core and optimal simultaneous protection}
\subsection{A simplex representation}
All weak $a$-players, meaning those at $L$ or $R$ labels, have identical
productivity. Replace the weak $a$-player in a tight diagonal triple by any
other weak $a$-player. The resulting inequality compares their payoffs;
reversing the replacement proves equality. Denote their common payoff by
$\alpha\ge0$. Write
\[
 l_s=p_{c_{L_s}}\ (0\le s\le q),\qquad
 r_s=p_{c_{R_s}},\ d_s=p_{b_{O_s}},\ o_s=p_{c_{O_s}}\ (1\le s\le q).
\]
Diagonal tightness gives
\[
 p_{b_{L_s}}=1-\alpha-l_s,\quad p_{b_{R_s}}=1-\alpha-r_s,\quad
 p_{a_{O_s}}=4-d_s-o_s.
\]
The critical triples $E_s,N_s$ respectively give
$o_s\ge l_{s-1}$ and $l_s\ge\alpha+r_s+d_s+o_s$.
Put $e_s=o_s-l_{s-1}$, $f_s=l_s-\alpha-r_s-d_s-o_s$, and
$h=p_{b_{L_q}}$. These variables satisfy
\begin{equation}\label{eq:simplex}
 (q+1)\alpha+l_0+\sum_{s=1}^q(d_s+r_s+e_s+f_s)+h=1,
 \qquad \alpha,l_0,d_s,r_s,e_s,f_s,h\ge0.
\end{equation}

\begin{proposition}\label{prop:simplex}
Equation~\eqref{eq:simplex}, with recursive reconstruction
\[
 o_s=l_{s-1}+e_s,\qquad
 l_s=l_{s-1}+\alpha+d_s+r_s+e_s+f_s,
\]
characterizes the entire core of $M_q$. In particular the core is affinely
isomorphic to a simplex of dimension $4q+2$.
\end{proposition}
\begin{proof}
Necessity follows from the displayed tight teams and inequalities.
For sufficiency, $l_s$ is nondecreasing and
$l_q=1-\alpha-h\le1-\alpha$, proving nonnegativity of all $b_L$ payoffs.
Also $r_s\le l_s-\alpha\le1-2\alpha$, giving nonnegativity of $b_R$.
Since $d_s+o_s\le l_s-\alpha-r_s\le1$, every strong $a$-payoff is at least
three. The other coordinates are visibly nonnegative and diagonal sums give
efficiency.

For a weak $a$-player, the positive $C$ edges of types
$(L_s,L_s),(R_s,R_s),(R_s,L_s),(L_{s-1},O_s)$ have total triple payoffs
\[
 1,\quad1,\quad1+l_s-r_s,\quad1+o_s-l_{s-1},
\]
respectively, all at least their worth one. If a strong $a_{O_s}$ has at
most one positive adjacent $A$ or $B$ entry, its own payoff of at least
three covers the triple's worth of at most three. When both are positive,
the four $(b,c)$ label pairs $(O_s,O_s),(O_s,L_s),(R_s,O_s),(R_s,L_s)$
have worths $4,4,4,5$ and respective payoff sums
\[
 4,\quad4+l_s-o_s,\quad5-\alpha-r_s-d_s,
 \quad5+l_s-\alpha-r_s-d_s-o_s.
\]
They cover their worths because $l_s\ge o_s$,
$\alpha+r_s+d_s\le l_s\le1$, and $f_s\ge0$. Zero-edge cases are
covered by nonnegativity. This exhausts all full triples and proves core
membership by~\eqref{eq:core}.

The parameters are recovered from the payoffs and differences already
specified, so reconstruction is injective. After replacing $\alpha$ by
$(q+1)\alpha$, equation~\eqref{eq:simplex} has $4q+3$ nonnegative
barycentric coordinates summing to one, hence dimension $4q+2$.
\end{proof}

\subsection{All ideals and the exact best fraction}
The simplex representation yields
\begin{equation}\label{eq:ideals}
 H_{a_{L_s}}=H_{a_{R_s}}=\frac1{q+1},\qquad H_{a_{O_s}}=4,
 \qquad H_{b_t}=H_{c_t}=1\quad\text{for every label }t.
\end{equation}
For the weak $a$ upper bounds use~\eqref{eq:simplex}; for strong $a$ use
diagonal tightness. The reconstructed formulas bound each $b,c$ coordinate
by one: in particular $o_s,d_s,r_s\le l_s\le1$.
Attainment uses separate simplex vertices. Set only $\alpha=1/(q+1)$ for
weak $a$ ideals; only $h=1$ for all strong $a$ and weak $b$ ideals;
only $d_s=1$ for $b_{O_s}$; only $r_s=1$ for $c_{R_s}$; and only $l_0=1$
for all $c_L,c_O$. No simultaneous attainment is assumed.

For any partial matching $P$ in any admissible market, let $U$ be its
unmatched players and $\delta=T-w(P)$. Every core point satisfies
$\sum_{i\in U}p_i\le\delta$. Consequently
\begin{equation}\label{eq:protectiondual}
 \rho(M)\le\frac{\delta}{\sum_{i\in U}H_i}
 \qquad\text{if }\sum_{i\in U}H_i>0.
\end{equation}
This follows directly by summing the desired floors over $U$; even lower
bounds on those ideals suffice for a weaker upper certificate.

For the triangular critical matching, $\delta=1$ and the unmatched set
in~\eqref{eq:unmatched} has $q+1$ weak $a$-players, each of ideal
$1/(q+1)$, and $2q+2$ players in sectors $b,c$, each of ideal one.
Thus $\sum_U H_i=2q+3$ and $\rho(M_q)\le1/(2q+3)$.
For equality put $t=1/(2q+3)$ and choose
\[
 \alpha=\frac{t}{q+1},\quad l_0=h=t,\quad
 d_s=r_s=t,\quad e_s=f_s=0.
\]
Equation~\eqref{eq:simplex} holds. Every $b,c$ payoff is at least $t$:
$l_s,o_s\ge t$, $d_s=r_s=t$, $b_{L_s}\ge h=t$, and
$b_{R_s}=1-\alpha-t\ge t$. Weak $a$-payoffs equal $tH_i$, and strong
$a$-payoffs are at least $3\ge4t=tH_i$. Hence
\[
                         \rho(M_q)=\frac1{2q+3}.
\]
Already $q=1$ gives a twelve-player market with best fraction $1/5<1/3$.
The optimizing point is strictly positive. A lottery over core allocations
cannot improve the simultaneous guarantee in expected payoffs: its mean
remains in the convex core. This observation makes no claim about other
criteria for evaluating lotteries.

\section{Selective subdivision and every cyclic allowed graph}
\subsection{A transfer lemma with explicit hypotheses}
Take a strict $k$-sector market with an optimal perfect matching relabeled
diagonally, worth $T$. Suppose edge $uv$ has matrix $B(i,j)$, with $i$ at $u$.
Choose $D_j\ge\max_i B(i,j)$. Insert a sector $d$, replacing $uv$ by
$ud$ with matrix $B$ and $dv$ with diagonal entries $D_j$ and off-diagonal
zeros. For an old team $t$ using $u$-label $i$ and $v$-label $j$, and a new
label $h$, its new worth is
\[
 w_1(t,h)=w_0(t)-B(i,j)+B(i,h)+D_h\mathbf1_{h=j}.
\]
For every old core point $z$, the unshifted lift $(z,D)$ covers every new
team. If $h=j$ this follows from the old inequality with the bonus added.
If $h\ne j$, use $D_h\ge B(i,h)$ and $B(i,j)\ge0$.
The extended diagonal matching and this cover have the same total
$T'=T+\sum_jD_j$, proving integral optimality and lifted core membership.
In particular all new diagonal teams are tight throughout the new core.

If $z$ has floor $\tau$, let $t=k\tau/(k+1)$. Subtract $t/k$ from every old
coordinate and give new-sector label $j$ payoff $D_j+t$. Every full-team
sum is unchanged. Equal sector sizes make the total change zero, and all
coordinates are at least $t$, so this shifted lift also lies in the core.

Now let an old partial matching $P$ have $r<m$ teams and worth $T-\delta$.
Extend each team using new label equal to its $v$-label. This remains
disjoint. If $P$ uses every $v$-label with positive $D_j$, it collects all
bonuses and keeps shortfall $\delta$. The new unmatched set has
$(k+1)(m-r)$ players. If the old floor witness has
$\tau=\delta/[k(m-r)]$, the shifted witness attains the new floor
$\delta/[(k+1)(m-r)]$. The unmatched budget proves that every unmatched
coordinate is fixed to that value on the entire first-stage optimal set.

An old unmatched target keeps ideal $H$ if its diagonal team has worth
$H$, its ideal is attained at an old core point, and the bonus at its
diagonal label is zero. The tight new diagonal gives the upper bound;
the unshifted lift of the attaining point gives the lower bound. More
generally every old ideal has at least its previous value by unshifted
lifting. No equality of entire cores or preservation of all ideals is needed.

\subsection{Applying the lemma repeatedly}
Start with the triangle's edge $B$ from $a$ to $c$. Its positive columns
are $\mathcal H=\{O_s,L_s:1\le s\le q\}$ and their maxima are two.
Set $D_j=2$ on $\mathcal H$ and zero otherwise. The critical matching
uses exactly these labels in sector $c$, so it collects every bonus.
Also $D_{L_0}=0$, preserving the target ideal one. Each subdivision adds
$4q$ to the total worth and keeps the unit shortfall.

Subdivide the edge still carrying $B$, oriented from $a$ to its current
neighbor, repeatedly. The matching labels at that neighbor remain exactly
$\mathcal H$, and all target diagonal labels remain $L_0$. Thus every
transfer hypothesis propagates. After $\ell-3$ steps the sector support is
an $\ell$-cycle, with
\[
 T_\ell=6q+1+4q(\ell-3),\quad w(P)=T_\ell-1,\quad
 |U|=\ell(q+1),\quad x_{c_{L_0}}=\frac1{\ell(q+1)},\quad H_{c_{L_0}}=1.
\]
All pair entries remain in $\{0,1,2\}$ and the shifted floor witness is
strictly positive.

Adding a zero-interaction sector leaves full-team worths unchanged.
An old core point lifts by zero; a floor witness is shifted exactly as
above with $D=0$. Extend the critical matching by any $r$ distinct new
labels. The same unit budget and first-stage pinning arguments apply;
the target's diagonal worth and attainable ideal stay one.

Choose a simple cycle in any fixed cyclic allowed graph $G$. Embed the
cycle construction there, add the other sectors with zero interactions,
and put zero matrices on all other allowed edges, including chords.
This is allowed by the model and gives the stated leximin target
$1/[K(q+1)]$ with ideal one.

For simultaneous protection retain the $2q+2$ original $b,c$ players in
$U$. Their old ideal-one witnesses lift without changing those payoffs;
they remain unmatched. Thus their enlarged-core ideals are at least one,
and~\eqref{eq:protectiondual} gives
$\rho(M_{G,q})\le1/(2q+2)$. We do not assert equality for longer cycles.
Letting $q$ grow gives $R(G)=0$, while nonnegativity and $c(G)\le R(G)$
give $c(G)=0$. Together with the tree proof, this proves
Theorem~\ref{thm:main}.

\section{Economic scope and reproducibility}
The classification identifies when stability itself permits a guarantee
relative to every participant's individually most favorable stable agreement.
For a tree, a single familiar fairness selection meets the sharp guarantee.
For a cyclic graph, the obstruction applies to the best possible selection.
Every finite market still has a positive optimal fraction; the zero guarantee
is a worst-case statement as the number of players per sector grows. The
strictly positive witness allocations separate this conclusion from an empty
core or a player who can never be paid positively.

The theorem concerns the entire labeled core and nominal leximin. It does not
cover productive partial-sector coalitions, unequal sector sizes, strictly
positive matrices on every allowed edge, ratio-normalized leximin, or a
selected proper face of the core. It does not determine the best constant at
each finite market size. These restrictions delimit the economic application,
rather than follow from the proof. In particular, an arbitrary transportation
or extreme-point representation cannot be assumed to preserve all of them.

\paragraph{Companion work.}
The protected-subsidy companion manuscript permits payments above grand-coalition
worth and asks how an external budget restores ideal-based floors. It shares
the strict pair-additive market class and individual-core-maxima benchmark,
but its sharp subsidy theorem uses a different extremal family. The present
paper's graph classification and exact within-core protection result concern
no external transfers. Earlier
chain--cycle and graph-theorem drafts are superseded presentations of this
paper, not additional contributions. The truthful-production companion studies
capacity reports and all-profile incentives, assumptions absent here.

\paragraph{Online Resource 1: reproducibility supplement.}
The self-contained supplement contains rational certificates, exact-solver
checks, instructions and source hashes. Tree checks compare cut coordinates
and clipping with raw full-team inequalities. Cyclic checks verify matching
and cover totals, shortfalls, floors and ideal attainment. The simultaneous
checks reconstruct simplex vertices and protection witnesses and test both
directions of the simplex description for specified finite parameters.
Negative controls reject damaged certificates. Exact coverage is listed in
the supplement. These computations corroborate the universal proofs above;
they are neither an exhaustive census of markets nor proof-assistant
formalization. No verification program is used as a premise of a theorem.

\section*{Statements and Declarations}
\paragraph{Funding.} No funding or support was received for this work.
\paragraph{Competing interests.} The author declares no competing interests.
\paragraph{Data and code availability.}
This theoretical study uses no observed participant or market data. The
constructed instances, source code and reproduction instructions are supplied
in Online Resource 1. A public archival identifier has not yet been assigned.
\paragraph{Use of artificial intelligence.}
OpenAI Codex assisted mathematical exploration, proof drafting, verification
code and manuscript preparation. Anthropic Claude was used for recorded
internal reviews during later preparation. The author personally reviewed
the manuscript, proofs, code and results. Internal model reviews and finite
checks are not human peer review or formal proof verification.

\begin{thebibliography}{99}
\small\setlength{\itemsep}{1pt}
\bibitem[Atay and N\'u\~nez(2019)]{AtayNunez2019}
Atay A, N\'u\~nez M (2019) Multi-sided assignment games on $m$-partite graphs.
Ann Oper Res 279:271--290. \url{https://doi.org/10.1007/s10479-019-03256-5}
\bibitem[Atay and Solymosi(2018)]{AtaySolymosi2018}
Atay A, Solymosi T (2018) On bargaining sets of supplier-firm-buyer games.
Econ Lett 167:99--103. \url{https://doi.org/10.1016/j.econlet.2018.03.022}
\bibitem[Gangam et~al.(2026)]{Gangam2026}
Gangam RR, Garg N, Shahkar P, Vazirani VV (2026) Equitable core imputations
for max-flow, MST and $b$-matching games. In: Proceedings of the 25th
International Conference on Autonomous Agents and Multiagent Systems,
IFAAMAS, pp 1619--1627. \url{https://doi.org/10.65109/VJNO3737}
\bibitem[Lucas(1995)]{Lucas1995}
Lucas WF (1995) Core theory for multiple-sided assignment games.
Duke Math J 81(1):55--65. \url{https://doi.org/10.1215/S0012-7094-95-08106-X}
\bibitem[Mart\'inez-de-Alb\'eniz et~al.(2020)]{Martinez2020}
Mart\'inez-de-Alb\'eniz FJ, Rafels C, Ybern N (2020) Assortative multisided
assignment games: The extreme core points. Games Econ Behav 120:144--153.
\url{https://doi.org/10.1016/j.geb.2020.01.001}
\bibitem[Mart\'inez-de-Alb\'eniz et~al.(2026)]{Martinez2026}
Mart\'inez-de-Alb\'eniz FJ, Rafels C, Ybern N (2026) Maximum and minimum core
payoffs for supermodular multisided assignment markets with local monotonicity.
Discrete Appl Math 379:605--612. \url{https://doi.org/10.1016/j.dam.2025.10.006}
\bibitem[N\'u\~nez and Rafels(2008)]{NunezRafels2008}
N\'u\~nez M, Rafels C (2008) A cooperative bargaining approach to the assignment
market. Group Decis Negot 17:553--563.
\url{https://doi.org/10.1007/s10726-008-9111-4}
\bibitem[Quint(1991)]{Quint1991}
Quint T (1991) The core of an $m$-sided assignment game.
Games Econ Behav 3(4):487--503. \url{https://doi.org/10.1016/0899-8256(91)90017-9}
\bibitem[Stuart(1997)]{Stuart1997}
Stuart HW Jr (1997) The supplier-firm-buyer game and its $m$-sided generalization.
Math Soc Sci 34(1):21--27. \url{https://doi.org/10.1016/S0165-4896(96)00830-X}
\end{thebibliography}
\end{document}
